\documentclass[preview]{standalone}

\usepackage{amsmath}
\usepackage{amssymb}
\usepackage{stellar}
\usepackage{definitions}
\usepackage{bettelini}

\begin{document}

\id{linearalgebra-notable-vector-spaces}
\genpage

\section{Transpose}

\includesnpt{matrix-transpose-definition}

\begin{snippetexample}{matrix-transpose-example}{Matrix transpose}
    Let
    \[
        A=
        \begin{pmatrix}
            2 & 3 & 7\\
            -1 & 8 & 2
        \end{pmatrix}
        \in\matrices_{2\times3}(\realnumbers).
    \]
    Then
    \[
        A^\transpose=
        \begin{pmatrix}
            2 & -1\\
            3 & 8\\
            7 & 2
        \end{pmatrix}
        \in\matrices_{3\times2}(\realnumbers).
    \]
\end{snippetexample}

\begin{snippetproposition}{matrix-transpose-linear-properties}{Transpose and linear operations}
    Let \(\mathbb{F}\) be a \field. If
    \(A,B\in\matrices_{m\times n}(\mathbb{F})\) and
    \(\lambda\in\mathbb{F}\), then
    \[
        (A+B)^\transpose=A^\transpose+B^\transpose,
        \qquad
        (\lambda A)^\transpose=\lambda A^\transpose.
    \]
\end{snippetproposition}

\begin{snippetproof}{matrix-transpose-linear-properties-proof}{matrix-transpose-linear-properties}{Transpose and linear operations}
    For every \(i,j\),
    \[
        ((A+B)^\transpose)_{ij}
        =
        (A+B)_{ji}
        =
        A_{ji}+B_{ji}
        =
        (A^\transpose+B^\transpose)_{ij}.
    \]
    The proof for scalar multiplication is analogous:
    \[
        ((\lambda A)^\transpose)_{ij}
        =
        (\lambda A)_{ji}
        =
        \lambda A_{ji}
        =
        (\lambda A^\transpose)_{ij}.
    \]
\end{snippetproof}

\section{Symmetric and skew-symmetric matrices}

\begin{snippetdefinition}{symmetric-skew-symmetric-matrix-definition}{Symmetric and skew-symmetric matrices}
    Let \(\mathbb{F}\) be a \field and let
    \(A\in\matrices_{n\times n}(\mathbb{F})\). The \matrix \(A\) is
    \emph{symmetric} if
    \[
        A^\transpose=A.
    \]
    The \matrix \(A\) is \emph{skew-symmetric} if
    \[
        A^\transpose=-A.
    \]
    We write
    \[
        \operatorname{Sym}(n)
        =
        \{A\in\matrices_{n\times n}(\mathbb{F})\suchthat A^\transpose=A\}
    \]
    and
    \[
        \operatorname{Skew}(n)
        =
        \{A\in\matrices_{n\times n}(\mathbb{F})\suchthat A^\transpose=-A\}.
    \]
\end{snippetdefinition}

\begin{snippetproposition}{symmetric-skew-symmetric-subspaces}{Symmetric and skew-symmetric matrix spaces}
    Let \(\mathbb{F}\) be a \field. The \set[sets]
    \[
        \operatorname{Sym}(n),
        \qquad
        \operatorname{Skew}(n)
    \]
    are linear subspaces of \(\matrices_{n\times n}(\mathbb{F})\).
\end{snippetproposition}

\begin{snippetproof}{symmetric-skew-symmetric-subspaces-proof}{symmetric-skew-symmetric-subspaces}{Symmetric and skew-symmetric matrix spaces}
    We prove the statement for \(\operatorname{Sym}(n)\). The zero \matrix is
    symmetric. If \(A,B\in\operatorname{Sym}(n)\), then
    \[
        (A+B)^\transpose=A^\transpose+B^\transpose=A+B,
    \]
    so \(A+B\in\operatorname{Sym}(n)\). If \(\lambda\in\mathbb{F}\), then
    \[
        (\lambda A)^\transpose=\lambda A^\transpose=\lambda A,
    \]
    so \(\lambda A\in\operatorname{Sym}(n)\).

    The skew-symmetric case is analogous. If \(A^\transpose=-A\) and
    \(B^\transpose=-B\), then
    \[
        (A+B)^\transpose=A^\transpose+B^\transpose=-A-B=-(A+B),
    \]
    and
    \[
        (\lambda A)^\transpose=\lambda A^\transpose=-\lambda A.
    \]
\end{snippetproof}

\begin{snippetproposition}{symmetric-skew-symmetric-decomposition}{Decomposition into symmetric and skew-symmetric parts}
    Let \(\mathbb{F}\) be a \field with
    \(\operatorname{char}(\mathbb{F})\neq2\). Then
    \[
        \matrices_{n\times n}(\mathbb{F})
        =
        \operatorname{Sym}(n)\oplus\operatorname{Skew}(n).
    \]
\end{snippetproposition}

\begin{snippetproof}{symmetric-skew-symmetric-decomposition-proof}{symmetric-skew-symmetric-decomposition}{Decomposition into symmetric and skew-symmetric parts}
    Let \(A\in\matrices_{n\times n}(\mathbb{F})\). Then
    \[
        A=
        \frac{A+A^\transpose}{2}
        +
        \frac{A-A^\transpose}{2}.
    \]
    The first \matrix is symmetric, because
    \[
        \left(\frac{A+A^\transpose}{2}\right)^\transpose
        =
        \frac{A^\transpose+A}{2}
        =
        \frac{A+A^\transpose}{2}.
    \]
    The second \matrix is skew-symmetric, because
    \[
        \left(\frac{A-A^\transpose}{2}\right)^\transpose
        =
        \frac{A^\transpose-A}{2}
        =
        -\frac{A-A^\transpose}{2}.
    \]
    Hence every \matrix is a sum of a symmetric \matrix and a skew-symmetric
    \matrix.

    If \(A\in\operatorname{Sym}(n)\intersection\operatorname{Skew}(n)\),
    then \(A^\transpose=A\) and \(A^\transpose=-A\). Hence \(A=-A\), so
    \(2A=0\). Since \(\operatorname{char}(\mathbb{F})\neq2\), \(A=0\).
    Therefore the intersection is \(\{0\}\), so the sum is direct.
\end{snippetproof}

\section{Dimensions}

\begin{snippetproposition}{symmetric-matrix-space-dimension}{Dimension of symmetric matrices}
    Let \(\mathbb{F}\) be a \field. Then
    \[
        \lineardim\operatorname{Sym}(n)=\frac{n(n+1)}{2}.
    \]
\end{snippetproposition}

\begin{snippetproof}{symmetric-matrix-space-dimension-proof}{symmetric-matrix-space-dimension}{Dimension of symmetric matrices}
    A symmetric \matrix is determined by its diagonal entries and by the
    entries above the diagonal; the entries below the diagonal are forced by
    symmetry. Thus the number of free parameters is
    \[
        n+(n-1)+\cdots+1=\frac{n(n+1)}{2}.
    \]
\end{snippetproof}

\begin{snippetproposition}{skew-symmetric-matrix-space-dimension}{Dimension of skew-symmetric matrices}
    Let \(\mathbb{F}\) be a \field with
    \(\operatorname{char}(\mathbb{F})\neq2\). Then
    \[
        \lineardim\operatorname{Skew}(n)=\frac{n(n-1)}{2}.
    \]
\end{snippetproposition}

\begin{snippetproof}{skew-symmetric-matrix-space-dimension-proof}{skew-symmetric-matrix-space-dimension}{Dimension of skew-symmetric matrices}
    If \(A=(a_{ij})\) is skew-symmetric, then
    \[
        a_{ii}=-a_{ii},
    \]
    so \(2a_{ii}=0\), and therefore \(a_{ii}=0\). Thus the diagonal is
    forced to be zero. The entries above the diagonal determine the entries
    below the diagonal, so the number of free parameters is
    \[
        (n-1)+(n-2)+\cdots+1=\frac{n(n-1)}{2}.
    \]
\end{snippetproof}

\begin{snippetexample}{symmetric-skew-symmetric-n3-example}{The case \(n=3\)}
    Let \(\mathbb{F}\) be a \field with
    \(\operatorname{char}(\mathbb{F})\neq2\). In
    \(\matrices_{3\times3}(\mathbb{F})\), a \basis of
    \(\operatorname{Sym}(3)\) is
    \[
        \{
            E_{11},E_{22},E_{33},
            E_{12}+E_{21},
            E_{13}+E_{31},
            E_{23}+E_{32}
        \},
    \]
    so \(\lineardim\operatorname{Sym}(3)=6\). A \basis of
    \(\operatorname{Skew}(3)\) is
    \[
        \{
            E_{12}-E_{21},
            E_{13}-E_{31},
            E_{23}-E_{32}
        \},
    \]
    so \(\lineardim\operatorname{Skew}(3)=3\). Hence
    \[
        6+3=9=\lineardim\matrices_{3\times3}(\mathbb{F}).
    \]
\end{snippetexample}

\section{An infinite-dimensional example}

\begin{snippetproposition}{real-numbers-over-rationals-infinite-dimensional}{\(\mathbb{R}\) over \(\mathbb{Q}\)}
    The \vectorspace \(\realnumbers\), viewed as a \vectorspace over
    \(\rationalnumbers\), is infinite-dimensional:
    \[
        \lineardim_{\rationalnumbers}\realnumbers=\infty.
    \]
\end{snippetproposition}

\begin{snippetproof}{real-numbers-over-rationals-infinite-dimensional-proof}{real-numbers-over-rationals-infinite-dimensional}{\(\mathbb{R}\) over \(\mathbb{Q}\)}
    It is enough to construct \linearlyindependent collections over
    \(\rationalnumbers\) of arbitrarily large finite size. Let
    \(p_1,\ldots,p_m\) be distinct prime numbers. We claim that
    \[
        \{\log p_1,\ldots,\log p_m\}
    \]
    is \linearlyindependent over \(\rationalnumbers\).

    Suppose
    \[
        \alpha_1\log p_1+\cdots+\alpha_m\log p_m=0,
        \qquad
        \alpha_i\in\rationalnumbers.
    \]
    Multiplying by a common denominator, we may assume that all coefficients
    are integers:
    \[
        a_1\log p_1+\cdots+a_m\log p_m=0,
        \qquad
        a_i\in\integers.
    \]
    Then
    \[
        \log(p_1^{a_1}\cdots p_m^{a_m})=0,
    \]
    so
    \[
        p_1^{a_1}\cdots p_m^{a_m}=1.
    \]
    By uniqueness of prime factorization, this is possible only if
    \(a_1=\cdots=a_m=0\). Hence also
    \(\alpha_1=\cdots=\alpha_m=0\).

    Since there are infinitely many primes, we can construct
    \linearlyindependent collections of arbitrary finite size. Therefore
    \(\realnumbers\) cannot have a finite generating collection as a
    \vectorspace over \(\rationalnumbers\).
\end{snippetproof}

\end{document}
